\scanpage{10}
\mspage{\handwritten{24}}

\margin{\handwritten{$/^{\circ}$ [\quad] inutile, résulte du \deleted{\unsure{fait|fait dit}} que $IB$ (\unsure{ou.|ou|resp.} $IC$) top. nilpotent.}}
\uline{contenant $\mathrm{IB}$ resp. $\mathrm{TC}$}\handwritten{$] /^{\circ}$}, enfin on suppose $\mathrm{IB}$ et $\mathrm{IC}$ fermés (condition automatiquement vérifiée si $\mathrm B,\mathrm C$ sont noethériens). \deleted{\unsure{Soit u:|Soit u}}Soient $\mathrm B_{\unsure{0|o}}=\mathrm B/\mathrm{IB}$,\deleted{et} $\mathrm C_0=\mathrm C/\mathrm{IC}$, $u:\inserted{\mathrm B\handwritten{\to}\mathrm C}$ un homomorphisme continu de $\mathrm A$-algèbres, $u_i:\mathrm B_i\handwritten{\to}\mathrm C_i$ l’homomorphisme \deleted{induit} déduit par passage au quotient (j’ai oublié d’introduire un système fondamental d’idéaux ouverts $\mathrm I_i\handwritten{\subset}\mathrm I_0\inserted{=\mathrm I}$ dans $\mathrm A$, et les $\mathrm B_i$ et $\mathrm C_i$ en conséquence\inserted{enfin ce sont tous les $\mathrm I_i\mathrm B$ et $\mathrm I_i\mathrm C$ qui doivent être fermés,}). Conditions équivalentes:

(i\handwritten{bis}) $u$ est \deleted{un isomorphisme} bijectif, et \tikz[baseline=(x.base)]{\node[draw,ellipse,inner sep=1pt] (x) {$u_0$ est un homéomorphisme.};}\handwritten{$\uparrow$}

\handwritten{$\curvearrowright$}(i\deleted{\unsure{bis}}) $u$ est un isomorphisme d’algèbres topologiques
% A curved arrow transposes this line to the position above (i bis).

(ii) Pour tout $i$, $u_i:\mathrm B_i\handwritten{\to}\mathrm C_i$ est bijectif, \tikz[baseline=(x.base)]{\node[draw,ellipse,inner sep=1pt] (x) {$u_0$ est un homéomorphisme.};}\handwritten{$\uparrow$}

(ii bis) pour tout $i$, muniss\deleted{\unsure{ons}}ant $\mathrm B_i$ et $\mathrm C_i$ des filtrations $\mathrm I$-adiques, l’homomorphisme $\operatorname{gr}_{\mathrm I}(\mathrm B_i)\handwritten{\to}\operatorname{gr}_{\mathrm I}(\mathrm C_i)$ induit par $u$\deleted{\unsure{e}} est bijectif, \tikz[baseline=(x.base)]{\node[draw,ellipse,inner sep=1pt] (x) {$u_0$ un homéomorphisme.};}\handwritten{$\uparrow$}
% slanted
\margin{\handwritten{$\bigcirc$ le \unsure{supposer} dès le début}}
% The ovals and upward pointers mark the repeated homeomorphism clauses.

Enfin, lorsque les $\deleted{\mathrm B}\handwritten{\mathrm C}_i$ sont plats sur les $\mathrm A_i$, \deleted{\unsure{les}}ces conditions équivalent également à

(iii) $u_0:\mathrm B_0\handwritten{\to}\mathrm C_0$ est un \deleted{isomorphisme\unsure{,}} \tikz[baseline=(x.base)]{\node[draw,ellipse,inner sep=1pt] (x) {homéomorphisme.};}\handwritten{$\uparrow$}

L’équivalence de (i) et (i bis) résulte trivialement du fait que les topologies de $\mathrm B,\mathrm C$ sont \deleted{adiques} définies par les puissances d’idéaux contenant $\mathrm{IB}$ resp. $\mathrm{IC}$. \deleted{Celle de} Evidemment (i\handwritten{bis})$\handwritten{\Rightarrow}$(ii), pour (ii)$\handwritten{\Rightarrow}$(i\handwritten{bis}) on note que les applications canoniques $\mathrm B\handwritten{\to}\underset{\handwritten{\longleftarrow}}{\lim}\mathrm B_i$ et $\mathrm C\handwritten{\to}\underset{\handwritten{\longleftarrow}}{\lim}\mathrm C_i$ sont bijectives grâce au fait que les $\mathrm I_i\mathrm B$ et $\mathrm I_i\mathrm C$ sont fermés (cf rectification à $0_{7.2.4.}$). L’équivalence de (ii) et (ii bis), \deleted{qu} mise pour mémoire, provient du fait que ($\mathrm I=\mathrm I_0$ étant un idéal de définition, donc nilpotent dans le\deleted{\unsure{s}}s $\mathrm A_i$) les filtrations mises sur $\mathrm B,\mathrm C$ sont finies. Enfin, reste à prouver que lorsque \deleted{\unsure{B est p}} $\deleted{\mathrm B}\handwritten{\mathrm C}_i$ est plat sur $\mathrm A_i$, alors $u_0$ bijectif implique $u_i$ bijectif ou si on veut $\operatorname{gr}(u)$ bijectif, ce qui est immédiat et sauf erreur figure déjà en lemme quelque part au Chap III (j’espère, avec la généralité qui convient, et bien entendu modulé ...). \handwritten{Cf \unsure{diagramme}}

\diagram{\handwritten{\begin{tikzcd}[ampersand replacement=\&]
\operatorname{gr} B \arrow[r,"\mathrm{surj}"] \& \operatorname{gr} C \\
B_0\otimes\operatorname{gr} A_i \arrow[u,"\mathrm{surj}"] \arrow[r,"\mathrm{bij}"'] \& C_0\otimes\operatorname{gr} A_{\unsure{i}} \arrow[u,"\mathrm{bij}"']
\end{tikzcd}}}
% archive sheet number: 9
